% Research continuation for the ProveIt repository, 11 October 2026.
% This file is assembled with the reviewed companion sections into the
% standalone deliverable Higher_Digamma_Powers_and_Ordered_Gamma_Sums.tex.
\documentclass[11pt,a4paper]{article}
\usepackage[T1]{fontenc}
\usepackage[utf8]{inputenc}
\usepackage{lmodern}
\usepackage[margin=27mm,headheight=15pt]{geometry}
\usepackage{amsmath,amssymb,amsthm,mathtools,mathrsfs}
\usepackage{booktabs,longtable,array}
\usepackage{microtype,xcolor,enumitem,fancyhdr}
\usepackage{xurl}
\usepackage[colorlinks=true,linkcolor=blue!45!black,citecolor=blue!45!black,urlcolor=blue!45!black]{hyperref}
\hypersetup{pdftitle={Higher Digamma Powers and Ordered Gamma Sums},
 pdfauthor={Research continuation prepared with ChatGPT for Vladimir Reshetnikov and ProveIt},
 pdfsubject={Exact finite parts, spectral zeta coefficients, and polylogarithmic depth cancellations}}
\numberwithin{equation}{section}
\newtheorem{theorem}{Theorem}[section]
\newtheorem{proposition}[theorem]{Proposition}
\newtheorem{lemma}[theorem]{Lemma}
\newtheorem{corollary}[theorem]{Corollary}
\theoremstyle{definition}
\newtheorem{definition}[theorem]{Definition}
\theoremstyle{remark}
\newtheorem{remark}[theorem]{Remark}
\DeclareMathOperator{\FP}{FP}
\DeclareMathOperator{\CT}{CT}
\DeclareMathOperator*{\Res}{Res}
\DeclareMathOperator{\Li}{Li}
\newcommand{\dd}{\mathrm d}
\newcommand{\Hrel}{\mathcal H}
\newcommand{\C}{\mathbb C}
\newcommand{\N}{\mathbb N}
\newcommand{\Q}{\mathbb Q}
\newcommand{\Z}{\mathbb Z}
\newcommand{\Log}{\operatorname{Log}}
\newcommand{\zet}{\zeta}
\newcommand{\commit}{b7667ae40a2957eed9a53d8559d53a3ee20dd1c1}
\setlength{\parskip}{0.35em}
\setlength{\emergencystretch}{2em}
\setlist{itemsep=3pt,topsep=5pt}
\allowdisplaybreaks[2]
\pagestyle{fancy}
\fancyhf{}
\fancyhead[L]{\small Higher Digamma Powers and Ordered Gamma Sums}
\fancyhead[R]{\small ProveIt research continuation}
\fancyfoot[C]{\thepage}

\begin{document}
\hypersetup{pageanchor=false}
\begin{titlepage}
\vspace*{14mm}
\noindent{\large\sffamily PROVEIT\quad /\quad RESEARCH CONTINUATION}
\vspace{18mm}

\noindent{\Huge\bfseries Higher Digamma Powers\\[3mm]
and Ordered Gamma Sums}
\vspace{8mm}

\noindent{\Large Exact finite parts, spectral zeta coefficients,\\[2mm]
and polylogarithmic depth cancellations}
\vspace{15mm}

\noindent Prepared with ChatGPT for Vladimir Reshetnikov and the ProveIt project\\
11 October 2026
\vspace{12mm}

\begin{abstract}
We develop a complete finite-part formula for every positive integer power
of the digamma function.  A normalized Mittag--Leffler expansion, finite
summation by parts, and a large-circle residue calculation reduce the
unit-interval moment to an absolutely convergent harmonic sum and the
linear Laurent coefficient of an explicit harmonic Dirichlet series.
The entire boundary polynomial is evaluated; its constant is zero for
odd powers and $(-1)^m\pi^{2m}/(2m+1)$ for power $2m$.  A second form
uses only an ordinary convergent series and an explicit finite list of
ordinary Stieltjes constants.  We carry out the fifth power completely,
including the mixed double-zeta derivatives that do not occur at powers
two through four, and give normalized primitives and differentiated
polygamma identities for all powers.

A complementary two-marker Gamma formula sums ordered relative harmonic
words over depth while retaining an arbitrary complex tagged exponent.
Integer marker differences give finite Hurwitz-zeta reductions, and
spectral differentiation produces generalized Stieltjes and log-Gamma
identities.  We also prove finite-interval formulas with several tagged
exponents and polylogarithmic depth cancellations.  Two precisely located
external-source errors are corrected by direct proofs.  Exact algebraic
checks and independent high-precision diagnostics accompany the proofs.
The Gaussian $S_6$ and revised $S_8$ conjectures remain unresolved.
\end{abstract}
\vfill
\noindent\small Source baseline: ProveIt revision
\texttt{\commit}.\\
The results are proposed additions to the inspected material; historical
priority beyond the cited comparison has not been established.
\end{titlepage}

\hypersetup{pageanchor=true}
\pagenumbering{roman}
\tableofcontents
\newpage
\pagenumbering{arabic}

\section{Research targets, scope, and notation}
\label{sec:scope}

The immediate starting point is the quartic digamma calculation in the
incoming reports \cite{Quartic,Barnes}.  The first higher-power question
in \cite{Quartic} asks for the fifth moment with \emph{every} finite
correction, including the correct linear Laurent coefficients and
ordinary zeta derivatives.  Its next question asks whether scalar
polynomials in the digamma function can cancel lower-depth data at
arbitrary degree.  The present article resolves the first question and
proves an all-power reduction mechanism.  It also identifies the precise
extra coordinate that must be addressed before the stronger scalar
cancellation proposal can be asserted at degree five.

The second line of work begins with the relative ordered harmonic
interpolation in the incoming independent-order report \cite{Independent}.
We prove a two-marker generating formula and exact depth cancellations
that keep a tagged complex exponent, together with all its spectral
derivatives.  This provides a useful connection between the harmonic
and polylogarithmic parts of the research program.

\begin{center}
\small
\begin{tabular}{>{\raggedright\arraybackslash}p{0.24\textwidth}
>{\raggedright\arraybackslash}p{0.39\textwidth}
>{\raggedright\arraybackslash}p{0.27\textwidth}}
\toprule
Source target & Result established here & Scope retained\\
\midrule
Fifth digamma finite part & Complete formula in Theorem~\ref{thm:q5},
with all scalar and spectral corrections & No reduction of the retained
spectral constants to a smaller arithmetic algebra\\[3pt]
Arbitrary integer powers & Normalized expansion, universal contour
constant, finite word algorithm, and ordinary subtracted sum & Fixed
integer power; no assertion for arbitrary complex powers\\[3pt]
Canonical scalar polynomials & Explicit fifth-degree cancellation of
height-one lower data; mixed derivative term remains & A formal
coordinate statement, not arithmetic independence\\[3pt]
Independent ordered exponents & Two-marker Gamma formula, integer-gap
reductions, Stieltjes and Gamma consequences & Stated convergence discs;
several tags are proved here on finite intervals\\[3pt]
Gaussian $S_6$, revised $S_8$ & A bounded expanded $S_6$ relation search
is recorded separately & Both remain conjectural\\
\bottomrule
\end{tabular}
\end{center}

The source review is targeted.  It is not a proof audit of every chapter
or every incoming archive.  The familiar Gamma recurrences and
asymptotics, finite stuffle algebra, and low-weight multiple-zeta
evaluations are antecedents, not new results \cite{DLMF,IKZ}.  Nor is the
existence of entire harmonic interpolation new: we use the established
interpolant of Ihara--Nakamura--Yamamoto \cite{INY}.  The new contribution
relative to the inspected repository is the collection of exact
extensions, normalizations, and reductions proved below.  No result is
accepted on the basis of integer-relation recognition.

\subsection{Conventions}

Let $\psi(z)=\Gamma'(z)/\Gamma(z)$, let $\gamma$ be Euler's constant,
and put
\[
 H_n^{(r)}=\sum_{j=1}^n j^{-r},\qquad H_n=H_n^{(1)},\qquad
 H_n(w)=\sum_{n\ge j_1>\cdots>j_d\ge1}
                     j_1^{-w_1}\cdots j_d^{-w_d}.
\]
The empty word has value $H_n(\varnothing)=1$; a nonempty word has
$H_0(w)=0$.  All finite words used in the digamma calculation have
positive integer entries.  We use the strict multiple-zeta convention
\begin{equation}
 Z(s_1,\ldots,s_d)=\zeta(s_1,\ldots,s_d)
 =\sum_{n_1>\cdots>n_d\ge1}
                 n_1^{-s_1}\cdots n_d^{-s_d}.
 \label{eq:strict-zeta}
\end{equation}
In particular $Z(s,w)=\sum_{n\ge1}H_{n-1}(w)n^{-s}$ where the series
converges.  At the positive integer words below, a first entry at least
two makes the series absolutely convergent.

Our Stieltjes convention is
\begin{equation}
 \zeta(1+u,a)=\frac1u+
       \sum_{j\ge0}\frac{(-1)^j}{j!}\gamma_j(a)u^j,
 \qquad a>0,\qquad \gamma_j=\gamma_j(1).
 \label{eq:stieltjes-convention}
\end{equation}
Thus $\gamma_0(a)=-\psi(a)$ and $\gamma_0=\gamma$.
For a meromorphic germ $f(u)$, $[u^j]f$ denotes its Laurent coefficient.
In particular $[u]f$ is meaningful even at a pole and is not then
$f'(0)$.  This distinction is essential in the fifth-power formula.

\begin{definition}[Original-coordinate finite part]
Suppose $f(x)=\sum_{j=1}^k a_jx^{-j}+g(x)$ near zero, with $g$
continuous on $[0,1]$.  Define
\begin{equation}
 \FP\int_0^1 f(x)\,\dd x
 =\lim_{\epsilon\downarrow0}\left\{
 \int_\epsilon^1 f(x)\,\dd x
 -\sum_{j=2}^k\frac{a_j\epsilon^{1-j}}{j-1}
 +a_1\log\epsilon\right\}.
 \label{eq:fp-definition}
\end{equation}
The coordinate is exactly $x$ and the upper endpoint is exactly $1$.
We write $Q_k=\FP\int_0^1\psi(x)^k\,\dd x$ for $k\ge1$.
\end{definition}

For a sequence $B_N=P(\log N)+o(1)$, with $P$ a polynomial, define
$\CT_{\log N}B_N=P(0)$.  This cutoff operation and a Laurent
coefficient in an auxiliary variable are not identified by convention;
Lemma~\ref{lem:cutoff} proves the comparison needed here.  Complex
logarithms in the global primitive use the principal branch on
$\C\setminus(-\infty,0]$.  All powers of positive real summation
variables use the real logarithm.

\section{Pole coefficients and the normalized global expansion}
\label{sec:poles}

The standard Taylor expansion of $\psi(1+\epsilon)$ and the recurrence
$\psi(z+1)=\psi(z)+1/z$ give \cite[\S\S5.5,5.7]{DLMF}
\begin{align}
 \psi(-n+\epsilon)
 &=-\frac1\epsilon+b_n+
                 \sum_{r\ge2}c_{r,n}\epsilon^{r-1},
 \label{eq:pole-series}\\
 b_n&=H_n-\gamma,\qquad
 c_{r,n}=H_n^{(r)}+(-1)^r\zeta(r),\qquad n\ge0.
 \notag
\end{align}
This is a convergent local expansion for $0<|\epsilon|<1$.

For each fixed positive integer $k$, define the finite polynomial data
\begin{align}
 p_{k,j}(n)
 &=[\epsilon^{k-j}]
   \left(-1+b_n\epsilon+\sum_{r=2}^k c_{r,n}\epsilon^r\right)^k,
      &&1\le j\le k, \label{eq:pdef}\\
 r_k(n)&=p_{k,1}(n),\qquad
 P_{k,n}(z)=\sum_{j=1}^k\frac{p_{k,j}(n)}{(z+n)^j},
 \qquad
 \kappa_k=[\epsilon^0]\psi(\epsilon)^k.
 \label{eq:Pdef}
\end{align}
The symbol $\kappa_k$ is a local Laurent constant; it will not be
confused with the universal contour constant $C_k$ introduced later.
The pole coefficients satisfy
\begin{equation}
 p_{k,j}(n)=O_k\bigl((1+\log(n+1))^{k-j}\bigr).
 \label{eq:p-growth}
\end{equation}
Indeed $b_n=O(1+\log(n+1))$ and every $c_{r,n}$ with $r\ge2$ is
bounded.  In a term of degree $k-j$ in $\epsilon$, at most $k-j$
factors of $b_n$ can occur.

\begin{theorem}[Normalized Mittag--Leffler formula for every power]
\label{thm:ML}
For $z\notin\{0,-1,-2,\ldots\}$,
\begin{equation}
 \boxed{\;
 \psi(z)^k=P_{k,0}(z)+\kappa_k+
       \sum_{n\ge1}\bigl(P_{k,n}(z)-P_{k,n}(0)\bigr).
 \;}
 \label{eq:ML}
\end{equation}
The series converges normally on compact sets avoiding the poles.
In particular the entire remainder is exactly the displayed constant,
with no uncomputed entire function.
\end{theorem}

\begin{proof}
Normal convergence of the candidate series follows from
\eqref{eq:p-growth}: on a fixed compact set, for sufficiently large $n$,
\[
 (z+n)^{-j}-n^{-j}=O(n^{-j-1}),
\]
uniformly in $z$.  It remains to determine the possible entire
remainder.  Put $f(w)=\psi(w)^k-P_{k,0}(w)$, so $f$ is analytic at
zero and $f(0)=\kappa_k$.  Integrate
\[
                  \frac{zf(w)}{w(w-z)}
\]
around $|w|=R_N=N+\tfrac12$, with positive orientation and $N$ large.

We need a uniform estimate on the full circle, including near the
negative real axis.  On its right half, the usual sectorial digamma
asymptotic gives $\psi(w)=O(1+\log R_N)$.  On its left half, use
\begin{equation}
             \psi(w)=\psi(1-w)-\pi\cot(\pi w).
 \label{eq:reflection}
\end{equation}
Here $1-w$ lies in the right half-plane and has size comparable to
$R_N$.  If $|\Im w|\ge1$, the cotangent is uniformly bounded.  If
$|\Im w|<1$, the points of this circle are a fixed positive distance
from the integers for all sufficiently large $N$, because its real
intersections are half-integers.  The cotangent is again uniformly
bounded.  Thus $f(w)=O_k((1+\log R_N)^k)$ on the whole contour, and
the integral is $O_k((1+\log R_N)^k/R_N)=o(1)$.

The residues at $w=z$ and $w=0$ are $f(z)$ and $-f(0)$.
At $w=-n$, the residue is
$-P_{k,n}(z)+P_{k,n}(0)$, as follows by expanding
$z/[w(w-z)]=1/(w-z)-1/w$ there.  The residue theorem and passage to
the limit give \eqref{eq:ML}.  This proves the normalization, not just
agreement of the principal parts.
\end{proof}

\subsection{A finite recurrence for the coefficient differences}

Write $\Delta p_{k,j}(n)=p_{k,j}(n)-p_{k,j}(n-1)$ for $n\ge1$,
and $\Delta r_k(n)=\Delta p_{k,1}(n)$.

\begin{proposition}[Exact difference recurrence]
\label{prop:difference}
For $1\le j\le k$,
\begin{equation}
 \Delta p_{k,j}(n)=
 \sum_{\ell=1}^{k-j}\binom{k}{\ell}
 \sum_{h=j}^{k-\ell}
       \binom{\ell+h-j-1}{h-j}
       \frac{p_{k-\ell,h}(n-1)}{n^{\ell+h-j}}.
 \label{eq:difference}
\end{equation}
An empty sum is zero; in particular $\Delta p_{k,k}(n)=0$.
\end{proposition}

\begin{proof}
The recurrence for $\psi$ gives the local identity
\[
 \psi(-n+\epsilon)=\psi(-(n-1)+\epsilon)+\frac1{n-\epsilon}.
\]
Expand its $k$th power by the binomial theorem.  In the term with
$\ell$ copies of $(n-\epsilon)^{-1}$, multiply the pole
$p_{k-\ell,h}(n-1)\epsilon^{-h}$ by
\[
 (n-\epsilon)^{-\ell}
   =\sum_{q\ge0}\binom{\ell+q-1}{q}
                       \frac{\epsilon^q}{n^{\ell+q}}.
\]
The coefficient of $\epsilon^{-j}$ has $q=h-j$, which gives
\eqref{eq:difference}.  Only finitely many coefficients contribute.
\end{proof}

The recurrence proves, in particular,
$\Delta p_{k,j}(n)=O_k((1+\log n)^k/n)$ for $j<k$.
Sharper logarithmic powers are immediate, but are not needed for the
exact identities.  Formula~\eqref{eq:difference} is also a convenient
implementation specification: it involves only lower powers and exact
rational binomial coefficients.

\section{The universal contour polynomial}
\label{sec:contour}

The subtraction constants in \eqref{eq:ML} are individually awkward
harmonic expressions.  Their combined cutoff asymptotic is much
simpler than their separate expansions suggest.

\begin{theorem}[All coefficients of the large-circle boundary term]
\label{thm:contour}
For every fixed integer $k\ge1$,
\begin{align}
 \kappa_k-\sum_{n=1}^N P_{k,n}(0)
 &=\frac1{2\pi}\int_{-\pi}^{\pi}
                         (\log N+i\theta)^k\,\dd\theta+o(1)
 \notag\\
 &=\sum_{m=0}^{\lfloor k/2\rfloor}
   \binom{k}{2m}\frac{(-1)^m\pi^{2m}}{2m+1}
                      (\log N)^{k-2m}+o(1).
 \label{eq:contour-polynomial}
\end{align}
Consequently
\begin{equation}
 C_k:=\CT_{\log N}\left(\kappa_k-\sum_{n=1}^N P_{k,n}(0)\right)
 =\begin{cases}
  (-1)^m\pi^{2m}/(2m+1),&k=2m,\\
  0,&k\text{ odd}.
 \end{cases}
 \label{eq:Ck}
\end{equation}
\end{theorem}

\begin{proof}
The residue at zero of $\psi(w)^k/w$ is $\kappa_k$.  At $-n$ it is
$-P_{k,n}(0)$.  Thus, on $R=N+\tfrac12$, the residue theorem gives
the exact identity
\begin{equation}
 \kappa_k-\sum_{n=1}^N P_{k,n}(0)
       =\frac1{2\pi}\int_{-\pi}^{\pi}
                           \psi(Re^{i\theta})^k\,\dd\theta.
 \label{eq:exact-circle}
\end{equation}
We compare this with the angular average of the principal logarithm.
On the right half-circle, $\psi(w)-\Log w=O(1/R)$ uniformly.
On the left half-circle with $|\Im w|\ge\log R$, reflection
\eqref{eq:reflection} and the exponential expansion of cotangent give
\[
 \psi(w)-\Log w=O(1/R)+O(e^{-2\pi|\Im w|}).
\]
For example, in the upper left quadrant,
$\Log(1-w)=\Log w-i\pi+O(1/R)$ and
$\pi\cot\pi w=-i\pi+O(e^{-2\pi\Im w})$; the lower quadrant is
the conjugate branch calculation.

The remaining arc near the negative real axis has angular length
$O(\log R/R)$.  On that arc, the bounded-cotangent argument in
Theorem~\ref{thm:ML} gives $\psi(w)-\Log w=O(1)$, while both
terms are $O(1+\log R)$.  Factoring the difference of the $k$th
powers shows that its contribution to the angular average is
$O_k((1+\log R)^k/R)$.  The rest of the circle contributes a quantity
of the same or smaller order.  The logarithm has a jump at the one
endpoint of the angular parametrization, which does not affect the
integral.  Therefore
\[
 \frac1{2\pi}\int_{-\pi}^{\pi}
   \bigl(\psi(Re^{i\theta})^k-(\log R+i\theta)^k\bigr)
                                      \,\dd\theta=o(1).
\]
Replacing $\log R$ by $\log N$ changes this by $o(1)$.  Odd powers
of $\theta$ integrate to zero and the even angular moments are
$(-1)^m\pi^{2m}/(2m+1)$, proving both assertions.
\end{proof}

For orientation, $C_1=C_3=C_5=0$, $C_2=-2\zeta(2)$, and
$C_4=18\zeta(4)$.  These are combined boundary constants.  They are
not the local Laurent constants $\kappa_k$.

\section{A finite-part formula for every positive integer power}
\label{sec:bridge}

Define
\begin{align}
 J_k&=\sum_{n\ge1}\sum_{j=2}^k
          \frac{\Delta p_{k,j}(n)}{(j-1)n^{j-1}},
 \label{eq:Jdef}\\
 D_k(s)&=\sum_{n\ge1}\frac{\Delta r_k(n)}{n^s},
                       \qquad \Re s>0.
 \label{eq:Ddef}
\end{align}
The first series is ordinary and absolutely convergent, by
Proposition~\ref{prop:difference}.  The second converges absolutely
in its initial half-plane.  Its meromorphic continuation to a
neighborhood of zero will follow from the proof below.

\begin{theorem}[The all-power harmonic spectral formula]
\label{thm:all-power}
For every integer $k\ge1$,
\begin{equation}
                 \boxed{\;Q_k=C_k+J_k+[s]D_k(s).\;}
 \label{eq:all-power}
\end{equation}
All data on the right are specified by the finite polynomial
calculation \eqref{eq:pdef}, the exact difference recurrence
\eqref{eq:difference}, and the coefficient of a meromorphic germ.
The identity uses the original-coordinate finite part
\eqref{eq:fp-definition}.
\end{theorem}

We first record the exact cutoff comparison that determines the sign
and the order of the spectral coefficient.

\begin{lemma}[Logarithmic cutoffs and Laurent coefficients]
\label{lem:cutoff}
Suppose
$a_n=P(\log n)+O(n^{-\delta}(1+\log n)^M)$, where $P$ is a
polynomial and $\delta>0$.  Let $A(s)=\sum_{n\ge1}a_nn^{-s}$
in its right half-plane of absolute convergence.  Then $A(1+u)$
continues meromorphically near $u=0$, and
\begin{align}
 \CT_{\log N}\sum_{n=1}^N\frac{a_n}{n}
           &=[u^0]A(1+u), \label{eq:cutoff-zero}\\
 \CT_{\log N}\sum_{n=1}^N\frac{a_n\log n}{n}
           &=-[u]A(1+u). \label{eq:cutoff-one}
\end{align}
\end{lemma}

\begin{proof}
For $P(L)=\sum_j c_jL^j$, write
\[
 A(1+u)=\sum_j c_j(-1)^j\zeta^{(j)}(1+u)
           +\sum_{n\ge1}\frac{a_n-P(\log n)}{n^{1+u}}.
\]
The last series and all its derivatives converge normally for
$\Re u>-\delta$.  The elementary Stieltjes limit formula gives
\[
 \sum_{n=1}^N\frac{(\log n)^j}{n}
       =\frac{(\log N)^{j+1}}{j+1}+\gamma_j
                    +O((1+\log N)^j/N).
\]
The stated error follows by one sum--integral comparison applied
to the derivative of $(\log x)^j/x$ on the tail.
On the Laurent side,
$[u^0](-1)^j\zeta^{(j)}(1+u)=\gamma_j$ and
$[u](-1)^j\zeta^{(j)}(1+u)=-\gamma_{j+1}$.
Termwise comparison proves the first formula and, after inserting
one further factor of $\log n$, the second.  Absolute convergence
handles the subtracted remainder in both cases.
\end{proof}

\begin{proof}[Proof of Theorem~\ref{thm:all-power}]
The harmonic-number asymptotic expansions, or reflection followed by
the positive-axis digamma expansion, show that
\begin{align}
 r_k(n)&=R_k(\log n)+O_k((1+\log n)^M/n),\label{eq:r-asymptotic}\\
 \Delta r_k(n)&=\frac{R_k'(\log n)}n
                  +O_k((1+\log n)^M/n^2),\label{eq:delta-asymptotic}\\
 R_k(L)&=\Res_{\epsilon=0}(L-\pi\cot\pi\epsilon)^k.
 \label{eq:Rdefinition}
\end{align}
Here and below $M$ may depend on $k$.  To justify the second estimate,
use the complete harmonic asymptotic expansions through the $1/n$
term: $r_k(n)=R_k(\log n)+n^{-1}U_k(\log n)
+O((1+\log n)^M/n^2)$ with a polynomial $U_k$.  The difference of
the explicit $n^{-1}U_k$ terms is $O((1+\log n)^M/n^2)$, and
the leading difference is $R_k'(\log n)/n$ to that same accuracy.
Thus the second formula does not follow from differencing an
unspecified first-order error.  It also follows directly from
\eqref{eq:difference} and the first two harmonic asymptotic terms.

Apply Lemma~\ref{lem:cutoff} to $a_n=n\Delta r_k(n)$.
It proves the meromorphic continuation of $D_k$ and gives
\begin{equation}
 \CT_{\log N}\sum_{n=1}^N\Delta r_k(n)\log n=-[s]D_k(s).
 \label{eq:delta-cutoff}
\end{equation}

Now integrate the normally convergent series \eqref{eq:ML}.
After the pole $P_{k,0}$ is removed, the same
$O(n^{-j-1})$ estimate proves uniform convergence also on the
closed interval $[0,1]$, so integration across its regular
endpoint is justified.
The finite part of the pole at zero contributes
$-\sum_{j=2}^kp_{k,j}(0)/(j-1)$; its simple pole has finite part
zero.  For a pole at $-n$, the integrated simple-pole term is
$r_k(n)\log(1+1/n)$, and the higher-pole terms have numerators
$p_{k,j}(n)(n^{1-j}-(n+1)^{1-j})/(j-1)$.
The constant subtraction integrates to $-P_{k,n}(0)$.

Use the following \emph{finite} summation identities:
\begin{align}
 &\sum_{n=1}^N p_{k,j}(n)
               \bigl(n^{1-j}-(n+1)^{1-j}\bigr)\notag\\
 &\hspace{6mm}=p_{k,j}(0)+
           \sum_{n=1}^N\frac{\Delta p_{k,j}(n)}{n^{j-1}}
                    -\frac{p_{k,j}(N)}{(N+1)^{j-1}},
                  &&j\ge2,\label{eq:sbp-high}\\
 &\sum_{n=1}^N r_k(n)\log(1+1/n)
     =r_k(N)\log(N+1)-\sum_{n=1}^N\Delta r_k(n)\log n.
 \label{eq:sbp-simple}
\end{align}
The first terms in \eqref{eq:sbp-high} cancel the finite part of
the pole at zero exactly, while its final terms tend to zero.
If $J_{k,N}$ denotes the cutoff in \eqref{eq:Jdef}, we obtain
\begin{align}
 Q_k=\lim_{N\to\infty}\biggl\{
 &\kappa_k-\sum_{n=1}^N P_{k,n}(0)+J_{k,N}
       +r_k(N)\log(N+1)\notag\\
 &\hspace{12mm}-\sum_{n=1}^N\Delta r_k(n)\log n\biggr\}.
 \label{eq:Qcutoff}
\end{align}
All terms here arose from finite sums or normally convergent series.
By \eqref{eq:r-asymptotic}, the polynomial asymptotic of the boundary
$r_k(N)\log(N+1)$ is $(\log N)R_k(\log N)$, which has no constant
term.  Theorem~\ref{thm:contour} supplies the constant $C_k$ for the
first two terms of \eqref{eq:Qcutoff}, and \eqref{eq:delta-cutoff}
supplies $+[s]D_k$.  Finally $J_{k,N}\to J_k$.  Taking the constant
coefficient of the complete cutoff expression proves
\eqref{eq:all-power}.
\end{proof}

\subsection{A closed ordinary-series form}

The residue polynomial in \eqref{eq:Rdefinition} itself has a closed
form.  This makes every Stieltjes correction explicit, without naming
any multiple-zeta Laurent coefficient.

% INSERT residue_polynomial.tex HERE

\begin{corollary}[Ordinary subtracted series at every power]
\label{cor:ordinary-all}
With the definitions above,
\begin{align}
 Q_k={}&C_k+J_k+
 k\sum_{m=0}^{\lfloor(k-2)/2\rfloor}
       (-1)^m\binom{k-1}{2m+1}\pi^{2m}\gamma_{k-2m-1}
 \notag\\
 &-\sum_{n\ge1}
   \left(\Delta r_k(n)-\frac{R_k'(\log n)}n\right)\log n.
 \label{eq:ordinary-all}
\end{align}
The finite Stieltjes sum is zero at $k=1$.  Both infinite sums are
ordinary absolutely convergent sums.  The last summand is
$O_k((1+\log n)^{M+1}/n^2)$.
\end{corollary}

This is useful both conceptually and computationally.  The right side
contains no unspecified regularization and no uncomputed entire
remainder.  Further terms of the same harmonic asymptotic expansions
give convergent Hurwitz-zeta tail accelerations if numerical evaluation
is wanted.  Such acceleration is optional for the identity itself.

\section{A finite algorithm in strict harmonic words}
\label{sec:algorithm}

The spectral formula becomes finite algebra once its target class is
specified.  Work over $\Q[\gamma,\zeta(2),\ldots,\zeta(k)]$ and put
$h_r=H_{n-1}^{(r)}$, $v=1/n$.  Compute the polynomial
$p_{k,j}$ from \eqref{eq:pdef}, and replace $H_n^{(r)}$ by
$h_r+v^r$ to obtain its value at $n$.  Subtract its value at $n-1$.
This produces every $\Delta p_{k,j}$ as a finite polynomial in
$v,h_1,\ldots,h_k$.

The finite stuffle recursion, with $(a,A)$ denoting a leading letter
followed by a word, is
\begin{equation}
 (a,A)*(b,B)=(a,A*(b,B))+(b,(a,A)*B)+(a+b,A*B).
 \label{eq:stuffle}
\end{equation}
It expresses products of the $h_r$ as a finite integer combination
of strict words $H_{n-1}(w)$.  The coefficients count equal-index
possibilities as well as the two strict orderings.  Thus
\[
        h_1^2=2H_{n-1}(1,1)+H_{n-1}^{(2)},
\]
and a further multiplication gives
\[
 h_1^3=6H_{n-1}(1,1,1)+3H_{n-1}(2,1)
                  +3H_{n-1}(1,2)+H_{n-1}^{(3)}.
\]
These are finite identities before any infinite sum is taken.

Every positive-integer word has a polynomial harmonic asymptotic
\begin{equation}
 H_{n-1}(w)=B_w(\log n)+O((1+\log n)^M/n)
 \label{eq:word-asymptotic}
\end{equation}
for some polynomial $B_w$ and integer $M$.  Here is a direct induction
that also supplies the continuation needed term by term.  For the
empty word the assertion is immediate.  Given it for $w$, use
$H_N(a,w)=\sum_{m=1}^N H_{m-1}(w)m^{-a}$.  If $a=1$,
the Stieltjes cutoff formula applied to each power of $\log m$
and the absolutely convergent subtracted remainder give a
polynomial in $\log N$ with an
$O((1+\log N)^{M'}/N)$ error.  If $a\ge2$, the sum converges
and its tail has that same error form, so the polynomial is a
constant.  Replacing $N$ by $n-1$ preserves the claimed form.
Consequently each $Z(1+s,w)$ has a meromorphic germ at zero by
Lemma~\ref{lem:cutoff}; for first entry at least two its defining
series is already holomorphic near zero.

\begin{theorem}[Terminating harmonic-word reduction]
\label{thm:algorithm}
At a fixed positive integer power $k$, the following procedure terminates
and yields the exact formula \eqref{eq:all-power}:
\begin{enumerate}
\item Compute $p_{k,j}$ and their finite differences by
\eqref{eq:pdef} or \eqref{eq:difference}.
\item Expand the summand of $J_k$ and the polynomial $\Delta r_k(n)$
in $v,h_1,\ldots,h_k$.
\item Apply the finite stuffle recursion to every harmonic monomial.
\item Replace $v^qH_{n-1}(w)$ by $Z(q,w)$ for $J_k$, and by
$Z(q+s,w)$ for $D_k(s)$.
\end{enumerate}
Every first entry in $J_k$ is at least two.  Every first entry in
$D_k(0)$ is at least one.  The resulting word weights are at most
$k-1$; coefficient factors in $\gamma$ and ordinary zeta values account
for the remaining formal weight.  The prescription introduces only a
finite number of specified linear Laurent coefficients.
\end{theorem}

\begin{proof}
Coefficient extraction and finite differences are finite polynomial
operations.  Each application of the stuffle recursion reduces the
combined remaining word length, so it terminates.  By
\eqref{eq:difference}, every monomial in $\Delta p_{k,j}$ contains
at least one $v$, and division by $n^{j-1}$ gives $q\ge j\ge2$
in $J_k$.  For $\Delta r_k$, it gives $q\ge1$.  Therefore all
replacements in the initial half-planes are identities of absolutely
convergent series.  Their continuations at zero agree by the
meromorphic subtraction justified by \eqref{eq:word-asymptotic}.
Finally assign weight $1$ to $h_1,\gamma,v$, weight $r$ to
$h_r,\zeta(r)$, and use \eqref{eq:pdef}.  Both relevant summand
polynomials have weight $k-1$, and stuffle preserves it.
\end{proof}

This is a reduction to a precise finite class of multiple-zeta values
and their spectral coefficients.  It does not prove that all those
values reduce further to products of ordinary zeta values.  For higher
powers it is appropriate to retain admissible multiple-zeta words
unless an exact additional relation is supplied.

\subsection{Recovery of the lower-power formulas}

Define the height-one coefficients
\begin{equation}
 \eta=[u]Z(1+u,1),\qquad
 \theta=[u]Z(1+u,1,1),\qquad
 \chi=[u]Z(1+u,1,1,1).
 \label{eq:height-coefficients}
\end{equation}
The first instances of Theorem~\ref{thm:all-power} are
\begin{align}
 Q_1={}&0,\notag\\
 Q_2={}&2\gamma_1-2\zeta(2),\label{eq:Q2}\\
 Q_3={}&3\zeta(2)-6\eta-6\gamma\gamma_1,\label{eq:Q3}\\
 Q_4={}&12\zeta(3)-12\gamma\zeta(2)+18\zeta(4)+4\zeta'(3)
       -24\theta+24\gamma\eta
                  +12(\gamma^2-\zeta(2))\gamma_1.
 \label{eq:Q4}
\end{align}
The cubic and quartic identities are consistency checks against the
incoming reports \cite{Quartic,Barnes}; they are not presented as new
evaluations here.  For example,
\begin{align*}
 D_2(s)&=-2\zeta(1+s),& J_2&=0,\\
 D_3(s)&=-6Z(1+s,1)+6\gamma\zeta(1+s),&J_3&=3\zeta(2),\\
 D_4(s)&=-24Z(1+s,1,1)+24\gamma Z(1+s,1)\\
       &\quad+12(\zeta(2)-\gamma^2)\zeta(1+s)+4\zeta(3+s),
       &J_4&=12\zeta(3)-12\gamma\zeta(2).
\end{align*}
At $k=1$ the residue is the constant $-1$, hence $J_1=D_1=0$;
the finite part also follows directly from $\log\Gamma(1)=0$ and
$\log\Gamma(\epsilon)=-\log\epsilon+O(\epsilon)$.

\section{The fifth power with every finite correction}
\label{sec:fifth}

For compactness write $\zeta_r=\zeta(r)$ in this section.  At a fixed
pole use $b=b_n$, $c=c_{2,n}$, $d=c_{3,n}$, and $e=c_{4,n}$.
Direct coefficient extraction gives
\begin{align}
 p_{5,5}&=-1,\qquad
 p_{5,4}=5b,\qquad
 p_{5,3}=5c-10b^2,\notag\\
 p_{5,2}&=10b^3-20bc+5d,\notag\\
 r_5&=-5b^4+30b^2c-10c^2-20bd+5e.
 \label{eq:p5}
\end{align}
For reference, the local constant is
\begin{equation}
 \kappa_5=-\gamma^5+20\gamma^3\zeta_2-55\gamma\zeta_4
             -30\gamma^2\zeta_3+20\zeta_2\zeta_3-5\zeta_5.
 \label{eq:kappa5}
\end{equation}
The universal constant is $C_5=0$, a different quantity.

For the finite differences, let $b,c,d$ instead denote the preceding
index $n-1$.  Then
\begin{align}
 \Delta p_{5,4}(n)&=\frac5n,\notag\\
 \Delta p_{5,3}(n)&=-\frac5{n^2}-\frac{20b}{n},\notag\\
 \Delta p_{5,2}(n)&=-\frac5{n^3}+\frac{10b}{n^2}
                              +\frac{30b^2-20c}{n},\notag\\
 \Delta r_5(n)&=\frac{-20b^3+60bc-20d}{n}
                          +\frac{10c}{n^2}+\frac{20b}{n^3}.
 \label{eq:delta5}
\end{align}
Equivalently the first numerator in the last line is $5r_4(n-1)$.
These identities can be checked simply by replacing
$(b,c,d,e)$ by $(b+1/n,c+1/n^2,d+1/n^3,e+1/n^4)$ in
\eqref{eq:p5}; Proposition~\ref{prop:difference} gives an independent
binomial derivation.

The higher poles therefore contribute
\begin{align}
 J_5
 &=\sum_{n\ge1}\left\{
       \frac{30(H_{n-1}-\gamma)^2
                    -20(H_{n-1}^{(2)}+\zeta_2)}{n^2}
                            -\frac{35}{6n^4}\right\}\notag\\
 &=30\gamma^2\zeta_2-60\gamma\zeta_3+\frac{35}{3}\zeta_4.
 \label{eq:J5}
\end{align}
For completeness, the low-weight evaluations needed in this last step
are
\[
 Z(2,1)=\zeta_3,\quad Z(2,1,1)=\zeta_4,\quad
 Z(2,2)=\frac34\zeta_4,\quad \zeta_2^2=\frac52\zeta_4.
\]
They follow from the standard duality, sum, and stuffle identities
\cite{DLMF,IKZ}.  In particular,
\begin{align*}
 \sum_{n\ge1}\frac{(H_{n-1}-\gamma)^2}{n^2}
       &=\frac{11}{4}\zeta_4-2\gamma\zeta_3+\gamma^2\zeta_2,\\
 \sum_{n\ge1}\frac{H_{n-1}^{(2)}+\zeta_2}{n^2}
       &=\frac{13}{4}\zeta_4,
\end{align*}
which proves every scalar term in \eqref{eq:J5}.

Applying finite stuffle to the last line of \eqref{eq:delta5} yields
the complete spectral identity
\begin{align}
 D_5(s)={}&-120Z(1+s,1,1,1)+120\gamma Z(1+s,1,1)
 \notag\\
 &+60(\zeta_2-\gamma^2)Z(1+s,1)
       +20(\gamma^3-3\gamma\zeta_2+\zeta_3)\zeta(1+s)
 \notag\\
 &+20Z(1+s,3)+10Z(2+s,2)+20Z(3+s,1)
 \notag\\
 &+10\zeta_2\zeta(2+s)-20\gamma\zeta(3+s).
 \label{eq:D5}
\end{align}
All equalities here hold first for $\Re s>0$, and then as meromorphic
germs.  In the first numerator of \eqref{eq:delta5}, the words
$(2,1)$ and $(1,2)$ from $H_{n-1}^3$ cancel against the corresponding
words in $H_{n-1}H_{n-1}^{(2)}$.  The residual $(3)$ word does not
cancel.  It is responsible for $Z(1+s,3)$.

Set
\begin{equation}
 D_{p,q}=[u]Z(p+u,q),\qquad
 \Omega_4=2D_{1,3}+D_{2,2}+2D_{3,1}.
 \label{eq:mixed-def}
\end{equation}
For $p\ge2$ these coefficients are ordinary outer-order derivatives,
$D_{p,q}=\partial_1Z(p,q)$.  For $p=1,q=3$, coefficient extraction
is required at a simple pole.

\begin{theorem}[Complete fifth-power finite part]
\label{thm:q5}
In the conventions \eqref{eq:fp-definition},
\eqref{eq:stieltjes-convention}, \eqref{eq:height-coefficients}, and
\eqref{eq:mixed-def},
\begin{align}
 Q_5={}&30\gamma^2\zeta_2-60\gamma\zeta_3+\frac{35}{3}\zeta_4
          -120\chi+120\gamma\theta
          +60(\zeta_2-\gamma^2)\eta
 \notag\\
 &-20(\gamma^3-3\gamma\zeta_2+\zeta_3)\gamma_1
          +20D_{1,3}+10D_{2,2}+20D_{3,1}
 \notag\\
 &+10\zeta_2\zeta'(2)-20\gamma\zeta'(3).
 \label{eq:Q5}
\end{align}
Equivalently the three mixed terms are $10\Omega_4$.
\end{theorem}

\begin{proof}
Use $C_5=0$, \eqref{eq:J5}, and the linear coefficient of
\eqref{eq:D5} in Theorem~\ref{thm:all-power}.  In particular
$[s]\zeta(1+s)=-\gamma_1$, which fixes the sign of its correction.
The other coefficients are exactly those named in
\eqref{eq:height-coefficients} and \eqref{eq:mixed-def}.
Every step rests on a finite algebraic identity or the normalized
analytic results already proved.
\end{proof}

This answers the fifth-power research question in \cite{Quartic}:
the fourfold height-one coefficient enters with multiplier $-120$,
and the mixed twofold coefficients and both ordinary zeta derivatives
are retained explicitly.  The result is not a reduction of $\chi$
or $\Omega_4$ to ordinary zeta and Stieltjes constants.

\subsection{An ordinary convergent formula for the same finite part}

The residue polynomial and its derivative are
\[
 R_5(L)=-5L^4+60\zeta_2L^2-90\zeta_4,
 \qquad R_5'(L)=-20L^3+120\zeta_2L.
\]
Consequently Corollary~\ref{cor:ordinary-all} gives the exact alternative
\begin{align}
 Q_5={}&30\gamma^2\zeta_2-60\gamma\zeta_3+\frac{35}{3}\zeta_4
                   +20\gamma_4-120\zeta_2\gamma_2
 \notag\\
 &-\sum_{n\ge1}
  \left\{\Delta r_5(n)
           +\frac{20\log^3n-120\zeta_2\log n}{n}\right\}\log n.
 \label{eq:Q5ordinary}
\end{align}
Here $\Delta r_5$ is the explicit harmonic expression
\eqref{eq:delta5}, with $b,c,d$ at $n-1$.
It can equally be written in polygamma values at positive integers:
\begin{equation}
 b=\psi(n),\qquad c=2\zeta_2-\psi'(n),\qquad
 d=\tfrac12\psi''(n).
 \label{eq:polygamma5}
\end{equation}
The final series in \eqref{eq:Q5ordinary} converges absolutely;
its summand is $O((1+\log n)^4/n^2)$.
This gives a definition-free numerical route to $Q_5$ using only
ordinary Stieltjes constants and ordinary polygamma evaluations,
with no multivariate analytic continuation routine.

\subsection{Explicit kernels for the height-one coefficients}
\label{subsec:heightkernel}

The retained height-one coefficient has a concrete normalized
integral representation.  We give it at every depth so that the
fourfold coefficient is specified independently of the digamma moment.
Height-one Laurent constants and related normalized integrals have
antecedents in Young's work \cite[Section~4]{Young}.  The proposition
below records a direct Mellin subtraction with the coefficient and
Stieltjes conventions fixed explicitly for this article.
Let
\begin{equation}
 \frac1{\Gamma(1+u)}=\sum_{j\ge0}a_ju^j,\qquad
 K_d(t)=\frac{[-\log(1-e^{-t})]^{d-1}}
                       {(d-1)!(e^t-1)},\quad d\ge1.
 \label{eq:Kd}
\end{equation}

\begin{proposition}[Normalized height-one Mellin coefficient]
\label{prop:heightkernel}
For $d\ge1$,
\begin{equation}
 Z(1+u,1^{d-1})=\sum_{j=0}^d a_ju^{j-d}+L_du+O(u^2),
 \label{eq:heightLaurent}
\end{equation}
where
\begin{align}
 L_d=a_{d+1}
 &+\int_0^1\log t\left\{K_d(t)
                  -\frac{(-\log t)^{d-1}}{(d-1)!t}\right\}\,\dd t
 \notag\\
 &+\int_1^\infty\log t\,K_d(t)\,\dd t.
 \label{eq:heightkernel}
\end{align}
Both integrals are absolutely convergent.  In particular
$L_2=\eta$, $L_3=\theta$, and $L_4=\chi$.
\end{proposition}

\begin{proof}
The generating series of $H_{n-1}(1^{d-1})$ gives, for $\Re u>0$,
\[
 \Gamma(1+u)Z(1+u,1^{d-1})=\int_0^\infty t^uK_d(t)\,\dd t.
\]
For example it follows by dividing
$\Li_{1^{d-1}}(z)=(-\log(1-z))^{d-1}/(d-1)!$
by $1-z$, multiplying by $z$, and setting $z=e^{-t}$; for $d=1$
the kernel is simply $1/(e^t-1)$.  Subtract the leading kernel on
$(0,1)$ to write this integral as $u^{-d}+B_d(u)$, with $B_d$
holomorphic near zero.  The subtraction is integrable with any fixed
power of $\log t$, because its remainder is at worst a polynomial in
$|\log t|$ near zero; the kernel decays exponentially at infinity.

The derivative identity
\[
 K_d(t)=-\frac1{d!}\frac{\dd}{\dd t}
                         [-\log(1-e^{-t})]^d
\]
shows that $B_d(0)=0$: the integrals from $\epsilon$ to infinity
and of the subtracted kernel from $\epsilon$ to $1$ have difference
\[
 \frac{[-\log(1-e^{-\epsilon})]^d-(-\log\epsilon)^d}{d!}
                         \longrightarrow0.
\]
Multiplication by $1/\Gamma(1+u)$ now proves
\eqref{eq:heightLaurent}, with $L_d=a_{d+1}+B_d'(0)$ exactly as
in \eqref{eq:heightkernel}.
\end{proof}

For the fifth moment the coefficient used in this kernel is explicitly
\begin{equation}
 a_5=\frac{\gamma^5-10\gamma^3\zeta_2+20\gamma^2\zeta_3
       +15\gamma\zeta_2^2-30\gamma\zeta_4
       -20\zeta_2\zeta_3+24\zeta_5}{120}.
 \label{eq:a5}
\end{equation}
This follows by exponentiating
$\log(1/\Gamma(1+u))=\gamma u+
\sum_{r\ge2}(-1)^{r+1}\zeta(r)u^r/r$.
The kernel is an exact analytic representation, not by itself a
reduction of the arithmetic content of $\chi$.

\subsection{The mixed derivative and its polylogarithmic representations}
\label{subsec:mixed}

The stuffle identity in a right half-plane gives
\[
 Z(1+u,3)=\zeta(1+u)\zeta_3-Z(3,1+u)-\zeta(4+u).
\]
Consequently
\begin{equation}
 D_{1,3}=-\gamma_1\zeta_3-\partial_2Z(3,1)-\zeta'(4).
 \label{eq:D13inner}
\end{equation}
The inner-order derivative has an absolutely convergent defining
series.  In particular this transformation replaces a singular
outer-order coefficient by a convergent inner-order derivative; it
does not remove the nonsymmetric spectral information.

Define the Tornheim function
\[
 T(a,b,c)=\sum_{m,n\ge1}\frac1{m^an^b(m+n)^c}
\]
in its absolute-convergence domain.  A finite partial fraction
calculation gives, for $u$ near zero,
\begin{equation}
 T(2,2,u)=2Z(2+u,2)+4Z(3+u,1).
 \label{eq:T22}
\end{equation}
Indeed, for $l=m+n$, summation over $1\le m<l$ uses
\[
 \frac1{m^2(l-m)^2}
 =\frac1{l^2}\left(\frac1{m^2}+\frac1{(l-m)^2}
                    +\frac2{m(l-m)}\right)
\]
and gives $2H_{l-1}^{(2)}/l^2+4H_{l-1}/l^3$.
Thus
\begin{align}
 \Omega_4
 &=2D_{1,3}+\tfrac12\partial_3T(2,2,0)
 \label{eq:OmegaT}\\
 &=2(\partial_1-\partial_2)Z(3,1)+\partial_1Z(2,2)
                      -2\gamma_1\zeta_3-2\zeta'(4).
 \label{eq:OmegaInner}
\end{align}
Spectral derivatives of Tornheim sums have substantial prior literature,
including \cite{BB2014}.  These equalities are useful coordinate changes;
no arithmetic independence is asserted for the displayed coordinates.

There are also completely convergent one-dimensional representations:
\begin{align}
 D_{1,3}
  &=-\gamma_1\zeta_3+
        \sum_{n\ge1}\frac{\zeta(3,n)\log n}{n},
 \label{eq:D13tail}\\
 \partial_3T(2,2,0)
  &=\int_0^\infty
        \frac{\Li_2(e^{-t})^2-e^{-t}\zeta_2^2}{t}\,\dd t.
 \label{eq:Tpolylog}
\end{align}
For \eqref{eq:D13tail}, use
$H_{n-1}^{(3)}=\zeta_3-\zeta(3,n)$ and extract the linear
coefficient.  For \eqref{eq:Tpolylog}, insert the Frullani identity
$-\log(m+n)=\int_0^\infty(e^{-(m+n)t}-e^{-t})\,\dd t/t$
and sum.  Absolute integrability follows either from
$\sum_{m,n\ge1}\log(m+n)/(m^2n^2)<\infty$, or directly from
the $O(\log t)$ integrand at zero and exponential decay at infinity.
Equations~\eqref{eq:OmegaT}, \eqref{eq:D13tail}, and
\eqref{eq:Tpolylog} connect the first mixed correction in a digamma
power to a dilogarithm integral and a Hurwitz-zeta tail sum.

\subsection{What scalar polynomial cancellation proves}
\label{subsec:canonical}

% INSERT canonical_fifth.tex HERE

\section{Normalized primitives and polygamma identities}
\label{sec:primitives}

The same normalization gives a primitive, not only a finite moment.
On $\C\setminus(-\infty,0]$ use the principal logarithm.

\begin{theorem}[An all-power primitive with fixed normalization]
\label{thm:primitive}
For $k\ge1$, put
\begin{align}
 \mathcal A_k(z)={}&\kappa_k z+p_{k,1}(0)\Log z
             +\sum_{j=2}^k\frac{p_{k,j}(0)z^{1-j}}{1-j}
 \notag\\
 &+\sum_{n\ge1}\biggl\{
       p_{k,1}(n)\log(1+z/n)-zP_{k,n}(0)
 \notag\\
 &\hspace{26mm}
       +\sum_{j=2}^k p_{k,j}(n)
             \frac{(z+n)^{1-j}-n^{1-j}}{1-j}\biggr\}.
 \label{eq:primitive}
\end{align}
The series converges normally on compact subsets of that slit plane,
and
\begin{equation}
             \mathcal A_k'(z)=\psi(z)^k,\qquad
             \mathcal A_k(1)=Q_k.
 \label{eq:primitive-properties}
\end{equation}
Its local Laurent--log expansion at zero has constant coefficient
zero, which fixes the otherwise arbitrary additive constant.
\end{theorem}

\begin{proof}
Each brace in \eqref{eq:primitive} is the integral from zero to $z$
of $P_{k,n}(w)-P_{k,n}(0)$.  For large $n$ and $z$ in a fixed
compact set its magnitude is
$O_k((1+\log n)^{k-1}/n^2)$, as follows by cancellation of its
linear terms or from the bound in the proof of Theorem~\ref{thm:ML}.
The normally convergent series may therefore be differentiated, and
\eqref{eq:ML} gives the first identity.  Every term in the braces
vanishes at zero and the analytic part begins with $\kappa_kz$.
The Laurent and logarithmic terms outside the braces also have no
constant coefficient.  The definition of the original-coordinate
finite part now gives $\mathcal A_k(1)=Q_k$.
\end{proof}

For $k=1$, this is the normalized branch of $\log\Gamma(z)$,
whose local expansion begins $-\Log z-\gamma z+\cdots$.
For every higher power it gives an explicitly normalized antiderivative
through harmonic pole coefficients.  Ordinary integrals along paths
contained in the slit plane follow by endpoint subtraction of
$\mathcal A_k$.  On other pole-free intervals the same conclusion
holds after choosing a consistent analytic continuation.

\begin{corollary}[Differentiated polygamma identities]
\label{cor:derivatives}
For $m\ge1$ and $z$ away from the nonpositive integers,
\begin{equation}
 \frac{\dd^m}{\dd z^m}\psi(z)^k
  =(-1)^m\sum_{n\ge0}\sum_{j=1}^k
              \frac{(j)_m p_{k,j}(n)}{(z+n)^{j+m}}.
 \label{eq:all-derivatives}
\end{equation}
The series is normally convergent.  Its left side is explicitly
\begin{equation}
 \sum_{m_1+\cdots+m_k=m}
          \frac{m!}{m_1!\cdots m_k!}
                         \prod_{i=1}^k\psi^{(m_i)}(z),
 \label{eq:polygamma-multinomial}
\end{equation}
where $\psi^{(0)}=\psi$ and $(j)_m$ is the rising factorial.
\end{corollary}

\begin{proof}
Differentiate \eqref{eq:ML} normally on compact subsets.  The constant
subtractions disappear after the first derivative and each pole gives
the displayed rising factorial.  The denominator exponent is at least
two, so \eqref{eq:p-growth} proves normal convergence directly as well.
The multinomial identity is the ordinary product rule.
\end{proof}

For example, with $b_n=H_n-\gamma$ the case $k=2,m=1$ gives
\[
 \psi(z)\psi'(z)
     =-\sum_{n\ge0}\left\{\frac1{(z+n)^3}
                             -\frac{b_n}{(z+n)^2}\right\}.
\]
This agrees with the classical polygamma pole series at $k=1$ and
extends it to polynomial products without an undetermined analytic
term.  The full formulas apply at arbitrary complex $z$ in their
specified domains, not only at integer arguments.

% INSERT ordered_section.tex HERE
% INSERT audit_section.tex HERE
% INSERT article_end.tex HERE
